\documentclass[12pt,a4paper]{article}

% ---- Encoding & Font ----
\usepackage{iftex}
\ifPDFTeX
  \usepackage[utf8]{inputenc}
  \usepackage[T1]{fontenc}
\else
  \usepackage{fontspec}  % XeTeX/LuaTeX (tectonic): Unicode nativo (·, —, ¿, ª, §)
\fi
\usepackage[spanish]{babel}
\usepackage{csquotes}

% ---- Page layout ----
\usepackage[top=2.5cm, bottom=2.5cm, left=2.5cm, right=2.5cm]{geometry}
\usepackage{setspace}
\onehalfspacing

% ---- Math ----
\usepackage{amsmath, amssymb, amsthm}
\usepackage{mathtools}
\usepackage{physics}
\usepackage{esint}  % \oiint

% ---- Graphics ----
\usepackage{graphicx}
\usepackage{tikz}
\usepackage{float}
\usepackage{caption}

% ---- Tables ----
\usepackage{array}
\usepackage{booktabs}
\usepackage{tabularx}
\usepackage{colortbl}

% ---- Colours ----
\usepackage{xcolor}
\definecolor{notebg}{HTML}{FFF3CD}
\definecolor{noteborder}{HTML}{FFC107}
\definecolor{boxred}{HTML}{F8D7DA}
\definecolor{boxredborder}{HTML}{F5C6CB}

% ---- Boxes ----
\usepackage{mdframed}
\usepackage{tcolorbox}
\tcbuselibrary{most}

\newtcolorbox{keybox}[1][]{
  colback=blue!5,
  colframe=blue!70!black,
  arc=4pt,
  boxrule=1pt,
  fonttitle=\bfseries,
  title={#1}
}

\newtcolorbox{warnbox}[1][]{
  colback=boxred,
  colframe=boxredborder,
  arc=4pt,
  boxrule=1pt,
  fonttitle=\bfseries,
  title={#1}
}

\newtcolorbox{notebox}[1][]{
  colback=notebg,
  colframe=noteborder,
  arc=4pt,
  boxrule=1pt,
  fonttitle=\bfseries,
  title={#1}
}

% ---- Hyperlinks & TOC ----
\usepackage[hidelinks]{hyperref}
\usepackage{bookmark}

% ---- Enumerate ----
\usepackage{enumitem}

% ---- Miscellaneous ----
\usepackage{icomma}
\usepackage{siunitx}
\usepackage{textcomp}
\usepackage[bottom]{footmisc}

% ---- Title ----
\title{\textbf{Clase 20: Materiales Magnéticos, Circuitos RL \\
        y Energía en el Campo Magnético}\\[0.3em]
        \large{Constante magnética, transitorio RL y densidad de energía}}
\author{Transcripto y expandido --- Curso Física III (OpenFING)}
\date{}

\begin{document}

\maketitle
\thispagestyle{empty}
\newpage

\tableofcontents
\newpage

% ===================================================================
\section{Materiales magnéticos en una bobina}
% ===================================================================

Como un dieléctrico en un condensador, un material magnético (núcleo de hierro)
cambia las propiedades de una bobina.

\begin{warnbox}[No lineales]
  A diferencia de los dieléctricos (lineales), la mayoría de los materiales
  magnéticos interesantes (hierro) son \textbf{altamente no lineales}. Esto
  permite la magnetización espontánea (imanes). No se trata aquí.
\end{warnbox}

Para materiales aproximadamente lineales:

\begin{keybox}[Constante magnética]
  \[
  \boxed{L = K_m\, L_{\text{vacío}}}
  \]
\end{keybox}

$L$ depende de la geometría de la bobina y del material.

% ===================================================================
\section{Circuitos RL}
% ===================================================================

\subsection{Planteo y ecuación diferencial}

Batería $\varepsilon$, resistencia $R$ y bobina $L$ en una malla, interruptor de
dos posiciones (A conecta, B desconecta). En $t=0$ se pone en A. Ley de mallas:
\[
\varepsilon - R I - L\frac{dI}{dt} = 0 \;\Longrightarrow\;
L\frac{dI}{dt} + R I = \varepsilon.
\]
Ecuación lineal de primer orden con segundo miembro.

\subsection{Analogía con el circuito RC}

Misma estructura que $\varepsilon = Q/C + R\,dQ/dt$: $R \leftrightarrow L$,
$1/C \leftrightarrow R$, $Q \leftrightarrow I$. Solución idéntica en forma.

\subsection{Resolución: solución particular y homogénea}

Particular: $I_0 = \varepsilon/R$ (estacionaria). Restándola, la homogénea
$L\,di/dt + R i = 0$ tiene solución exponencial $i = i_0 e^{rt}$ con
$r = -R/L$. Imponiendo $I(0)=0$:

\begin{keybox}[Corriente al conectar (posición A)]
  \[
  \boxed{I(t) = \frac{\varepsilon}{R}\left(1 - e^{-Rt/L}\right)}
  \]
\end{keybox}

\begin{warnbox}[Test físico]
  Para $t\to\infty$, $I\to\varepsilon/R$. Una exponencial creciente indicaría un
  error de cuentas.
\end{warnbox}

\subsection{Conexión (A) y desconexión (B)}

Al pasar a B (sin batería), la corriente decae:
\[
i(t) = i_0\, e^{-Rt/L}.
\]

\subsection{Tiempo característico}

\begin{keybox}[Constante de tiempo RL]
  \[
  \boxed{\tau = \frac{L}{R}}
  \]
\end{keybox}

Bobina grande $\to$ $\tau$ grande (tarda en almacenar energía); resistencia
grande $\to$ $\tau$ chico (domina la disipación).

% ===================================================================
\section{Experimentos de inducción}
% ===================================================================

Tres experimentos (ley de Faraday); el flujo varía por módulo, área o ángulo.

\begin{itemize}[nosep]
  \item \textbf{Ángulo:} imán girando en una bobina enciende una lamparita
        (generador); alimentado con corriente, funciona como motorcito.
  \item \textbf{Módulo:} anillo de aluminio sobre una bobina con CA y núcleo de
        hierro levita (anillo de Thomson); cortado, no levita; sin núcleo,
        tampoco.
  \item \textbf{Corrientes parásitas:} péndulo de bronce macizo que se frena
        entre bobinas encendidas; con cortes (peine), oscila libre. Principio de
        los núcleos laminados de transformadores.
\end{itemize}

% ===================================================================
\section{Energía almacenada en una bobina}
% ===================================================================

\subsection{Balance de potencia}

Con $\varepsilon = R I + L\,dI/dt$, la potencia entregada:
\[
P = \varepsilon I = R I^2 + L I \frac{dI}{dt}.
\]
$R I^2$ es la disipada; el resto es la energía almacenada por unidad de tiempo:
$dU_B/dt = L I\,dI/dt$.

\subsection{Energía magnética}

Con $I\,dI/dt = \tfrac12 \,d(I^2)/dt$:

\begin{keybox}[Energía en una bobina]
  \[
  \boxed{U_B = \frac{1}{2}L\,I^2}
  \]
\end{keybox}

Constante elegida para $U_B=0$ si $I=0$. Análoga a $U_E = \tfrac12 Q^2/C$.

% ===================================================================
\section{Densidad de energía del campo magnético}
% ===================================================================

Tomando el solenoide ideal ($B = \mu_0 n I$, sección $A$, largo $l$, volumen
$lA$, $L = \mu_0 N^2 A/l$) y dividiendo la energía por el volumen, con
$I = Bl/(\mu_0 N)$:

\begin{keybox}[Densidad de energía magnética]
  \[
  \boxed{u_B = \frac{B^2}{2\mu_0}}
  \]
\end{keybox}

Probada para el solenoide, vale en general (en el vacío).

\begin{notebox}[Analogía]
  $u_E = \tfrac12 \varepsilon_0 E^2$ y $u_B = \tfrac{1}{2\mu_0}B^2$: ambas
  proporcionales al cuadrado del campo; $\mu_0$ aparece en el denominador por
  convención histórica de Biot--Savart.
\end{notebox}

\begin{notebox}[Próxima clase]
  Cable coaxial, circuito LC, analogía masa--resorte y resonancia.
\end{notebox}

\end{document}
